\documentclass[11pt,letterpaper]{article}

% ── Geometry ──────────────────────────────────────────────────────────────────
\usepackage[margin=1in]{geometry}

% ── Pagination ────────────────────────────────────────────────────────────────
\usepackage[all]{nowidow}    % Require ≥2 lines on each side of a page break
\brokenpenalty=10000         % No hyphenated words across page breaks
\raggedbottom                % Allow uneven page bottoms rather than stretching

% ── Fonts ─────────────────────────────────────────────────────────────────────
\usepackage[T1]{fontenc}
\usepackage{newpxtext}     % Palatino-based serif (TeX Gyre Pagella)
\usepackage{newpxmath}     % Matching math font

% ── Colors ────────────────────────────────────────────────────────────────────
\usepackage[dvipsnames]{xcolor}
\definecolor{SectionBlue}{HTML}{1B3A5C}
\definecolor{AccentGray}{HTML}{4A4A4A}
\definecolor{QuoteBorder}{HTML}{8B9DAF}
\definecolor{RiskRed}{HTML}{C0392B}
\definecolor{RiskAmber}{HTML}{E67E22}
\definecolor{RiskGreen}{HTML}{27AE60}
\definecolor{BoxBg}{HTML}{F7F8FA}

% ── Tables ────────────────────────────────────────────────────────────────────
\usepackage{booktabs}
\usepackage{tabularx}

% ── Boxes ─────────────────────────────────────────────────────────────────────
\usepackage[framemethod=tikz]{mdframed}

% ── Lists ─────────────────────────────────────────────────────────────────────
\usepackage{enumitem}
\usepackage{needspace}

% ── Hyperlinks ────────────────────────────────────────────────────────────────
\usepackage{hyperref}
\hypersetup{
  colorlinks=true,
  linkcolor=SectionBlue,
  urlcolor=SectionBlue,
  citecolor=SectionBlue,
  pdfauthor={Punt Labs},
  pdftitle={Lux PR/FAQ: A Visual Output Surface for Terminal-Hosted AI Coding Agents},
  pdfsubject={Working Backwards: Hypothesis Stage},
}

% ── Bibliography ─────────────────────────────────────────────────────────────
\usepackage[style=numeric,backend=biber,sorting=none]{biblatex}
\addbibresource{prfaq.bib}

% ── Section styling ───────────────────────────────────────────────────────────
\usepackage{titlesec}
\titleformat{\section}
  {\needspace{6\baselineskip}\Large\bfseries\color{SectionBlue}}
  {}
  {0pt}
  {}

\titleformat{\subsection}
  {\needspace{5\baselineskip}\large\bfseries\color{AccentGray}}
  {}
  {0pt}
  {}

\titleformat{\subsubsection}
  {\normalsize\bfseries\color{AccentGray}}
  {}
  {0pt}
  {}
\titlespacing*{\subsubsection}{0pt}{1em}{0.5em}[0pt]

% ── Document stage ───────────────────────────────────────────────────────────
\newcommand{\prfaqstage}[1]{\def\theprfaqstage{#1}}
\prfaqstage{hypothesis}

% ── Document version ─────────────────────────────────────────────────────────
\newcommand{\prfaqversion}[2]{\def\theprfaqmajor{#1}\def\theprfaqminor{#2}}
\prfaqversion{1}{3}

% ── Header/footer ────────────────────────────────────────────────────────────
\usepackage{fancyhdr}
\pagestyle{fancy}
\fancyhf{}
\fancyhead[L]{\footnotesize\textcolor{AccentGray}{Stage: \theprfaqstage{} \enspace\textbar\enspace v\theprfaqmajor.\theprfaqminor}}
\fancyhead[R]{\footnotesize\textcolor{AccentGray}{\thepage}}
\renewcommand{\headrulewidth}{0pt}

% ══════════════════════════════════════════════════════════════════════════════
% Custom environments
% ══════════════════════════════════════════════════════════════════════════════

\newcommand{\prsection}[1]{%
  \vspace{1em}%
  {\large\bfseries\color{SectionBlue}#1}%
  \vspace{0.3em}\par%
}

\newmdenv[
  topline=false,
  bottomline=false,
  rightline=false,
  linewidth=3pt,
  linecolor=QuoteBorder,
  backgroundcolor=BoxBg,
  innerleftmargin=12pt,
  innerrightmargin=12pt,
  innertopmargin=8pt,
  innerbottommargin=8pt,
  skipabove=\baselineskip,
  skipbelow=\baselineskip,
]{customerquote}

\newmdenv[
  topline=false,
  bottomline=false,
  rightline=false,
  linewidth=3pt,
  linecolor=SectionBlue,
  backgroundcolor=BoxBg,
  innerleftmargin=12pt,
  innerrightmargin=12pt,
  innertopmargin=8pt,
  innerbottommargin=8pt,
  skipabove=\baselineskip,
  skipbelow=\baselineskip,
]{spokespersonquote}

\newcounter{faqnum}
\newenvironment{faqpair}[1]{%
  \needspace{4\baselineskip}%
  \vspace{0.5em}%
  \refstepcounter{faqnum}%
  \noindent\textbf{\color{SectionBlue}Q\thefaqnum: #1}\par\nopagebreak[4]%
  \vspace{0.2em}%
  \begin{adjustwidth}{1em}{0pt}%
  \setlength{\parindent}{0pt}%
}{%
  \end{adjustwidth}%
  \vspace{0.5em}%
}
\newcommand{\faqref}[1]{\hyperref[#1]{FAQ~\ref*{#1}}}

\newcounter{featurenum}
\newcommand{\featureitem}[2]{%
  \refstepcounter{featurenum}%
  \item[\textbf{\color{SectionBlue}F\thefeaturenum.}]%
  \textbf{#1}: #2%
}
\newcommand{\featureref}[1]{\hyperref[#1]{Feature~\ref*{#1}}}

\usepackage{changepage}

% ══════════════════════════════════════════════════════════════════════════════
% Document
% ══════════════════════════════════════════════════════════════════════════════

\begin{document}

% ── Headline block ───────────────────────────────────────────────────────────
\begin{center}
  {\LARGE\bfseries\color{SectionBlue}
    Lux Gives Terminal-Based AI Coding Agents a Visual Canvas} \\[0.8em]
  {\large\color{AccentGray}
    For the agent: describe UI as a typed JSON element tree, and Lux renders it
    natively with GPU-accelerated Dear~ImGui on the developer's own machine, no
    vendor compute} \\[0.3em]
  {\large\color{AccentGray}
    For the developer: see complex agent output as filterable tables,
    dashboards, and diagrams instead of scrollback, and drive a terminal-based
    agent through the GUI. A fixed, curated element set that composes without
    limit}
\end{center}

\vspace{0.5em}

% ══════════════════════════════════════════════════════════════════════════════
% PRESS RELEASE
% ══════════════════════════════════════════════════════════════════════════════

\section*{Press Release}

AUSTIN, TX, January 2027. Punt Labs today released Lux, a local-first
visual output surface for terminal-hosted AI coding agents. The most capable
agentic coding tools, Claude Code\cite{anthropic2025claudecode} and
Codex\cite{openai2025codex}, run in the terminal on the developer's own
machine, but their structured output is trapped in scrollback: a table, a
dependency diagram, or a metrics dashboard arrives as monospace text or an
ad~hoc screenshot. Lux gives those agents a screen. An agent describes what it
wants as a tree of typed JSON elements; Lux renders it with Dear~ImGui\cite{imgui2025}
in a persistent, GPU-accelerated window, entirely on the developer's machine,
with no vendor cloud. Unlike web dashboards that require running the agent
on vendor compute, or the terminal that has no structure at all, Lux keeps the
agent local and gives its output a surface the developer can read, filter,
and click.

\prsection{Problem}

The most capable coding agents, Claude Code and Codex, live in the terminal:
full tool use, MCP integration, multi-turn autonomous operation, and output that
is all text. That power comes with no graphical surface. When an agent produces a
table of failing tests, a diagram of a module graph, or a dashboard of build
status, the developer reads it as raw scrollback. There is no way to filter the
table, select a row, or keep the view live as the underlying data changes. The
most powerful agents run exactly where there are no graphics.

A second problem sits on top of the first: developers rarely run just one agent.
Several run at once, across repositories, sessions, and machines, and tracking
what each is doing is hands-on and error-prone. This is an aggregation problem,
many independent agents and one developer with no shared place to watch them.

Vendors are answering with hosted web dashboards that often require running
the \emph{entire agent} (not just the model call) on the vendor's compute: the
developer trades local control for a viewport. IDE-hosted tools and MCP-host
apps solve pieces of this but not the whole problem: one manages agents with
text-only output, the other renders rich interfaces but cannot drive a
terminal agent's autonomous, multi-turn workflow end-to-end
(see~\faqref{faq:competition}). The result is a split: the tools with the most
agentic power have no visual output, and the tools with visual output lack the
agentic power.

\prsection{Solution}

Lux is a persistent local service that agents write to over MCP. An agent calls
a tool such as \texttt{show\_table} or \texttt{show\_dashboard} with structured
content; the display renders it in a native GPU-accelerated window using
Dear~ImGui. Multiple agents share the one window, each in its own frame with
clear attribution, so a developer glances at a single surface and sees what
every session is doing. Everything runs on the developer's machine: the Hub
holds authoritative UI state, the display renders a replica, and no data
leaves the box.

Lux is not a passive canvas. It ships two capabilities that matter in daily
work. First, a rich data-display core: filterable, searchable tables with row
selection and detail panels; dashboards that compose metric cards, charts, and
tables; diagrams and images. The element vocabulary is deliberately limited: a
curated, fixed set of built-in kinds that compose without limit, the same way
a small set of ImGui widgets does. Second, and the direction Lux is heading, an
ask-the-user interaction loop: an agent can pose a question through a dialog,
checkbox, or button and receive the answer back. When the user clicks, the
interaction routes back to the owning Hub, the real handler runs there, and the
Hub re-sends the updated view. That loop is the first interaction primitive, the
seed of a two-way surface where the user drives a terminal-based agent through a
GUI instead of only reading its output. The canonical example is the beads issue
browser: a filterable table that auto-refreshes on every \texttt{bd} command
through a post-tool hook, running live in every active agent session
(see~\faqref{faq:evidence}).

\prsection{Customer Quote}

\begin{customerquote}
\textit{``I keep four Claude Code sessions open across three repos. Before Lux I
had four terminals and no idea what any of them were doing without scrolling
back. Now one window shows me a beads board, a test-run table, and a build
dashboard, one frame per agent, all live. When a migration script needed me
to confirm a destructive step, the agent didn't dump a wall of text and wait;
it put a dialog on the screen with the two options. I clicked. That is the
difference between watching a transcript and using a tool.''}

\raggedleft --- \textbf{Maya Torres}, Staff Engineer at a Series~B DevTools Startup
\end{customerquote}

\prsection{Getting Started}

Getting started takes one command:

\begin{enumerate}[nosep]
  \item Install:

    {\small\texttt{curl -fsSL https://raw.githubusercontent.com/\allowbreak punt-labs/lux/main/install.sh | sh}}
  \item Restart Claude Code. The Lux display window opens automatically the
    first time any agent sends visual output.
  \item Ask an agent to show you something: ``show me the beads board,''
    ``put the test results in a table.'' The view appears in the window, live.
\end{enumerate}

No API key is required. Lux runs on macOS and Linux, needs Python~3.13+, and
installs the \texttt{punt-lux} package plus the Claude Code plugin.

\prsection{Spokesperson Quote}

\begin{spokespersonquote}
``The most powerful coding agents live in the terminal, and the terminal gives
them one output mode: text. We built Lux so an agent can hand the developer a
real surface (a table you filter, a dashboard you scan, a dialog you click)
without moving the agent onto someone's cloud. The part I care most about is
the interaction loop: the agent asks, the user answers, and the answer flows
back to the handler that asked. That turns a display into a conversation. And
because the whole thing is local and native, `Lux is showing me stale data' is
not a failure mode: the Hub owns the state and re-sends the view when it
changes.''

\raggedleft --- \textbf{Jim Freeman}, Punt Labs
\end{spokespersonquote}

\prsection{Call to Action}

Lux is free and open-source under the MIT license. Visit
{\small\texttt{https://github.com/punt-labs/lux}} to install, read the target
architecture, or contribute a new element kind.

\newpage

% ══════════════════════════════════════════════════════════════════════════════
% FREQUENTLY ASKED QUESTIONS
% ══════════════════════════════════════════════════════════════════════════════

\section*{Frequently Asked Questions}

% ── External FAQs (Customer-facing) ──────────────────────────────────────────

\subsection*{External FAQs}

\begin{faqpair}{What is Lux and who is it for?}\label{faq:what}
  Lux is a visual output surface for terminal-hosted AI coding agents. It is a
  persistent, GPU-accelerated window, built on Dear~ImGui\cite{imgui2025},
  that agents write to and read from. It is for developers who run terminal
  agents such as Claude Code\cite{anthropic2025claudecode} or
  Codex\cite{openai2025codex} and want their structured output as something they
  can see and act on: filterable tables, dashboards, diagrams, and interactive
  prompts, rather than scrollback text. The element kinds are fixed and
  curated; they compose without limit, covering any layout the way a small
  set of ImGui widgets does. Everything renders and runs on the developer's
  own machine. Lux runs on macOS and Linux.
\end{faqpair}

\begin{faqpair}{How is this different from terminal scrollback, Claude Desktop, or vendor dashboards?}\label{faq:competition}
  Terminal agents (Claude Code, Codex) are the most capable but output only
  text. Claude Desktop renders rich interfaces and is itself an MCP host, but it
  cannot drive a coding session end-to-end: it lacks the terminal agent's
  autonomous, filesystem-level, multi-turn tool-driven workflow. Vendor web
  dashboards provide visuals but often require running the whole agent on
  vendor compute. IDEs such as Cursor manage agents with text-only output.

  Lux fills the gap: a native visual surface for terminal-hosted agents, running
  on the developer's machine. It is not a chat interface and not a cloud
  console. It is a screen the agents share: one window, one frame per agent,
  live. See \faqref{faq:competitors} for the detailed competitive analysis.
\end{faqpair}

\begin{faqpair}{How do I get started?}
  One command installs Lux and registers the Claude Code plugin. Restart Claude
  Code. The display window opens by itself the first time an agent sends visual
  output. There is no separate app to launch and no configuration to edit. Ask
  an agent to show you a table or a board and it appears.
\end{faqpair}

\begin{faqpair}{How much does it cost?}
  Free and open-source under the MIT license. There is no paid tier, no account,
  and no API key requirement for the display itself. Lux consumes the same agent
  you already run; it adds no per-call cost of its own.
\end{faqpair}

\begin{faqpair}{What happens to my data?}
  Everything stays local. Agents talk to the Hub over MCP; the Hub talks to the
  display over a Unix domain socket under \texttt{/tmp/lux-<user>/}. Per-project
  configuration lives in your repository at
  \texttt{<repo>/.punt-labs/lux.md}. There is no telemetry and no cloud
  dependency: nothing about your code, data, or agents' output leaves the
  machine. The Hub owns authoritative UI state and the display renders
  a replica; render calls never cross the wire, only UI state does.
\end{faqpair}

\begin{faqpair}{What can Lux render today, and can agents ask me questions?}\label{faq:capabilities}
  Two things ship and are in daily use. The data-display core covers filterable
  and searchable tables with row selection and detail panels, dashboards
  (metric cards, charts, tables), diagrams, and images: 25 element kinds in
  all. The interaction loop lets an agent ask you a question through the UI: a
  dialog, a checkbox, or a button, with your answer routed back to the agent's
  handler. That ask-the-user loop shipped with the Hub/Display interaction model.
  Completing coverage of every ImGui widget as a built-in element is the v2
  direction (see~\faqref{faq:roadmap}).
\end{faqpair}

% ── Internal FAQs (Business-facing) ──────────────────────────────────────────

\newpage
\subsection*{Internal FAQs}

\subsubsection*{Value \& Market}

\begin{faqpair}{What is the total addressable market?}\label{faq:tam}
  The primary market is developers running terminal-hosted AI coding agents:
  Claude Code\cite{anthropic2025claudecode} and Codex are the leading tools,
  both terminal-hosted. That population is a subset of professional developers,
  expanding as agentic coding moves from novelty to daily practice. We decline
  a precise headcount: no reliable public figure for active terminal-agent
  developers exists to cite without inventing it. As an order-of-magnitude
  \emph{assumption} only (not a measurement), the cohort is plausibly
  $10^5$--$10^6$ today, a small fraction of the world's professional
  developers, and growing.

  Positioning, not sizing, is the point. Lux is free and MIT-licensed; its value
  to Punt Labs is ecosystem pull, not license revenue (see~\faqref{faq:revenue}).
  The relevant question is not ``how many will pay'' but ``how many terminal-agent
  developers would keep a local visual surface open once they have one.'' That is
  the hypothesis current internal use supports and external use has not yet
  tested.
\end{faqpair}

\begin{faqpair}{What evidence do we have that customers want this?}\label{faq:evidence}
  Lux is released (v0.19.1 on PyPI as \texttt{punt-lux} and on the plugin
  marketplace) and used daily. The precise scope matters and we state it
  plainly: the daily user is \emph{one} primary human developer (Jim Freeman)
  plus the agent fleet he runs across the Punt Labs repositories. When we say the
  beads issue browser (a filterable table that auto-refreshes on every
  \texttt{bd} command through a post-tool hook) ``runs in every session,''
  that counts \emph{agent} sessions, not distinct human users. It is a genuine
  daily-driver signal, but n${=}$1 human. That single view still demonstrates the
  two claims the product rests on: the data-display core has real daily value,
  and the live-refresh hook makes the surface feel current without the agent
  having to coordinate updates. The interaction loop, an agent asking a
  question and receiving the answer through a dialog, is exercised in real
  sessions, not just tests.

  The honest gap: this evidence is \emph{internal}. The target customer (an
  external developer running several terminal agents at once, the profile in the
  press-release quote) has not used Lux yet. Punt Labs is the current
  \emph{user} and proof point, not the market. External adoption is zero: no
  usage data, no retention figures, no third-party testimonials. Validation
  confined to the company that built the product is, in the four-risks framing, a
  High Value-risk signal, and this document rates it as such. Every ``value''
  claim here should be read against that boundary (see~\faqref{faq:next-step} and
  the Value risk in the Risk Assessment).
\end{faqpair}

\begin{faqpair}{Who are the competitors and why will we win?}\label{faq:competitors}
  The landscape divides on two axes: agentic power and visual surface.

  \begin{itemize}[nosep]
    \item \textbf{Claude Code / Codex.} Terminal-hosted, the most capable
      agentic tools: full tool use, MCP, multi-turn autonomy. No visual output
      beyond scrollback. \emph{These are Lux's integration targets, not rivals.}
    \item \textbf{Claude Desktop.} Renders novel visual interfaces and is the
      canonical MCP host, but lacks the terminal agent's autonomous,
      filesystem-level, multi-turn workflow end-to-end.
    \item \textbf{Cursor.} IDE with strong agent management, but structured
      output stays text. Extensible via VS~Code extensions, which are IDE-scoped
      and hand-written.
    \item \textbf{Vendor web dashboards.} Provide visual management, but often
      require running the whole agent on vendor compute. The developer gives up
      local control for a viewport.
    \item \textbf{Jupyter.} Rich visual output via the display protocol, but
      notebook-first, not agent-first.
  \end{itemize}

  \textbf{Primary threat: Anthropic-native visual output.} Anthropic has every
  incentive to build visual panels into Claude Code. A basic built-in display
  would narrow Lux's advantage to ``native GPU-accelerated and fully local vs.
  vendor-native.'' Survival rests on three structural facts a competitor cannot
  copy cheaply: Lux renders natively with ImGui on the developer's GPU and keeps
  the agent local; Lux is a shared surface across \emph{multiple} agents and
  tools, not a panel bolted onto one vendor's client; and the Hub/Display split
  lets an agent on a remote box drive a GUI on your local screen, because UI
  state rather than render calls crosses the wire (see~\faqref{faq:architecture}).
  A web-native competitor would have to abandon web-only rendering to match the
  local-native path: a platform decision they are unlikely to reverse.
\end{faqpair}

\begin{faqpair}{What is the next step to validate external demand?}\label{faq:next-step}
  The riskiest assumption is that Lux's internal value transfers to developers
  outside Punt Labs. The first external-validation experiment is to put Lux in
  front of external developers who run 2+ concurrent agent sessions daily and
  measure whether they keep the window open past the first day. Concretely:
  recruit a cohort, ship them the one-command install, and instrument
  nothing on their machine (no telemetry by design); instead, interview at
  day~7. The
  decision threshold: if fewer than half describe the display as something they
  would miss if it were removed, revisit the positioning before investing in the
  v2 element-coverage buildout.
\end{faqpair}

\begin{faqpair}{How do external developers discover an MIT-licensed tool?}\label{faq:acquisition}
  Three concrete channels, none of them yet proven for Lux.

  \begin{itemize}[nosep]
    \item \textbf{The Claude Code plugin marketplace.} Lux is published there; a
      developer browsing plugins can find and install it. This is the most direct
      channel and requires no separate discovery of Punt Labs.
    \item \textbf{GitHub.} The repository is public and MIT-licensed; the
      one-command install lives in the README. This is organic search and
      word-of-mouth: real, but slow and unmeasured.
    \item \textbf{Cross-surfacing from the companion suite.} A developer who
      adopts vox, biff, or quarry encounters Lux as the surface those tools
      render into, and vice versa. The suite is the distribution flywheel the
      ecosystem thesis depends on (see~\faqref{faq:revenue}).
  \end{itemize}

  The honest position: distribution is currently a hope, not a proven motion.
  ``Organic and word-of-mouth'' is not a strategy, and we have no acquisition
  data. The marketplace listing and the suite cross-surfacing are the two levers
  we can pull; whether either converts external installs is exactly what
  the day-7 cohort (see~\faqref{faq:next-step}) begins to answer.
\end{faqpair}

\subsubsection*{Technical}

\begin{faqpair}{What is the architecture?}\label{faq:architecture}
  Lux is a three-tier distributed Hub/Display system. The target design lives
  in {\small\texttt{docs/architecture/\allowbreak target/\allowbreak target.md}}.

  \begin{itemize}[nosep]
    \item \textbf{\texttt{luxd}}: the session Hub. It decodes wire elements,
      holds authoritative UI state, ownership, and handler dispatch in
      \texttt{HubDisplay}, pushes element replicas to the display, and
      re-dispatches interactions that come back from clicks.
    \item \textbf{\texttt{lux-display}}: the ImGui renderer. It stores a full
      copy of the UI, renders it every frame, and wraps handlers so that a click
      forwards back to the owning Hub instead of executing locally.
    \item \textbf{\texttt{mcp-proxy}}: the transport bridge from Claude
      Code's stdio to the \texttt{luxd} WebSocket.
  \end{itemize}

  MCP is a first-class front door into the Hub, not a thin adapter. The
  boundary rule is strict: UI state crosses the wire (scene
  trees, serialized element objects, interaction messages); render calls do not.
  Interactive elements use two-tier handler dispatch (the Hub owns the real
  handler and the Display forwards the interaction back), so the Hub and
  Display can be separate processes, potentially on separate machines.

  This lets the three tiers sit on different machines. The Display connects to
  the Hub over a Unix socket; agents connect to the Hub over a WebSocket that
  binds to localhost (127.0.0.1:8430) by default. A remote agent works today,
  two ways: tunnel that localhost port over SSH, or start the Hub with
  {\small\texttt{luxd --host}} pointed at a non-local interface for a direct TCP
  connection (a session key authenticates either way). The practical payoff:
  your most powerful agent can run on a remote box and still drive a GUI on your
  local screen, because only UI state crosses the wire.
\end{faqpair}

\begin{faqpair}{What makes the engineering distinctive?}\label{faq:formal}
  The display-process concurrency is formally specified in Z and model-checked
  with ProB. The model-check covers only the display singleton
  spawn/bind/reap \emph{lifecycle}, not the rendering, protocol, or
  handler-dispatch surface: it is not a claim that the whole system is
  formally verified. Multiple agents may simultaneously decide no display is
  running and race to become the one process that owns a single Unix-socket
  path. Exactly one must win. That spawn/reap/bind singleton lifecycle is
  modeled in \texttt{docs/display\_lifecycle.tex} as a finite state machine over
  concurrent agents, type-checked \texttt{fuzz}-clean, and exhaustively
  model-checked. ProB enumerates every interleaving and confirms five
  invariants: singleton-serving (never two live owners), never-unlink-live (a
  probe cannot delete a socket whose owner is alive), no-two-winners across the
  bind window, lost-racer-clean-exit (the process that loses the race exits
  before opening a window), and deadlock-freedom under the two-lock spawn/bind
  discipline.

  This matters because the failure it prevents is exactly the one that bit us in
  testing: a second agent probing during the window between \texttt{bind} and
  \texttt{listen} reads the fresh socket as dead, unlinks it, and opens a second
  window. Proving the interleaving space exhaustively, rather than testing a
  few orderings, is rare rigor for a developer tool, and the specification is
  a regression artifact for any future lifecycle change.
\end{faqpair}

\begin{faqpair}{What is the v1-to-v2 roadmap, and what is the migration risk?}\label{faq:roadmap}
  \textbf{v1 (shipped):} the agent shows the developer complex data and
  visualizations. The rich data-display core, plus the ask-the-user loop
  (dialog, checkbox, button) as the first interaction primitive, on the
  Hub/Display model.

  \textbf{v2 (direction):} the developer interacts with the agent through a GUI.
  The surface turns two-way, so the user drives a terminal-based agent through a GUI
  instead of only viewing its output. Completing ImGui widget coverage is the
  means to that end, not the end in itself: richer interaction needs more input
  primitives (sliders, inputs, selects) and the composites that arrange them.
  This is still Punt Labs building one fixed, curated vocabulary, explicitly
  \emph{not} agents authoring arbitrary extensions and not a plugin registry,
  because a bounded element set already composes without limit.

  The honest risk is the migration itself. Lux is mid-rewrite from an older
  single-process design onto a distributed Element-ABC path with two-tier remote
  handler dispatch. Four of the 25 element kinds (text, button, checkbox, and
  dialog) are on the new Element-ABC path today; the other 21 are legacy
  dataclasses being migrated family by family into curated built-ins. The
  sequencing, per-kind contract, and open design decisions are mapped in
  {\small\texttt{docs/architecture/\allowbreak element-migration-audit.md}}. The scope is real work
  (interactive inputs need a handler-wrapping seam, composites need typed
  children, the table kind carries built-in view state), and it is the primary
  feasibility risk (see the Risk Assessment). It is mitigated by incremental
  crossing: each kind flips independently and mixed scenes keep working, with the
  legacy path deleted only after the last kind crosses.
\end{faqpair}

\begin{faqpair}{What is the development approach and delivery order?}\label{faq:timeline}
  Delivery is ordered by dependency, not by calendar. The pattern is established
  on the lowest-risk element kinds first and then climbs interactivity and
  composite complexity:

  \begin{enumerate}[nosep]
    \item Reconcile the four migrated exemplars into one source of truth.
    \item Invert the handler-wrapping seam so each interactive element declares
      its own interaction (prerequisite for migrating the input kinds).
    \item Basics display-only leaves (image, separator, progress, spinner,
      markdown).
    \item Interactive value inputs (slider, combo, text/number inputs, radio,
      color picker, selectable).
    \item Simple composites (group, tab bar, collapsing header).
    \item Stateful composites (window, modal).
    \item Draw and plot.
    \item Table, on its own (it carries built-in filter/selection state).
    \item Retire the legacy render path once every kind has crossed.
  \end{enumerate}

  What gets cut if scope compresses: v2 element coverage is a fast-follow, not a
  v1 blocker. The data-display core and the interaction loop already ship; the
  migration improves internal structure and widens the built-in vocabulary
  without gating the value already in daily use.
\end{faqpair}

\begin{faqpair}{If Lux succeeds, what breaks first?}\label{faq:scaling}
  Three known limits, stated with their current status rather than reassurance.

  \begin{enumerate}[nosep]
    \item \textbf{The image texture cache grows without eviction.} Every image
      an agent shows uploads a texture that is never freed: an unbounded
      memory leak under sustained image use. Status: known and tracked only as a
      Should-Do (a texture-cache eviction policy, \featureref{feat:texture}); it
      is a real limitation today, not a solved problem. It is the first thing
      that breaks under a long-running session that renders many distinct images.
    \item \textbf{Renderer maintainability debt.} \texttt{display/server.py} is
      roughly 1,550 lines, well over the 300-line target, and
      \texttt{element\_renderer.py} is also over. Decomposition is in progress
      (the original monolith was already split once). This does not break at
      runtime, but it slows every rendering change and is why the render layer is
      undertested. Status: actively being decomposed, incrementally.
    \item \textbf{Many-agent contention on one window.} The shared window is a
      single surface; a large agent fleet all pushing scenes contends for one
      display and one Hub. The whole-tree resend replication policy is simple but
      re-sends the whole affected UI on change. Status: adequate at the current
      fleet size; a heavy multi-agent load is untested and would be the first
      place the simple replication model needs a diff protocol.
  \end{enumerate}
\end{faqpair}

\subsubsection*{Business}

\begin{faqpair}{What is the revenue model?}\label{faq:revenue}
  Lux has no direct revenue. It is free and MIT-licensed. Its strategic value is
  ecosystem pull for the Punt Labs agent-tool suite. Lux is the visual
  counterpart to the other tools in the suite: vox (voice), biff (presence), and
  quarry (search). It surfaces them, too: the beads board, presence panels, and
  status views render in the Lux window. A developer who adopts Lux is one step
  from the rest of the suite. On the cost side there is no COGS (Lux runs
  entirely on the developer's own machine), so Punt Labs' development and
  maintenance time is the only real expense. This is an infrastructure and
  ecosystem play, acceptable because it is treated as a public good with
  strategic return, not a profit center. If Lux does not pull adoption of the companion
  tools, the investment is a public good without business return: that is
  the viability risk.
\end{faqpair}

\begin{faqpair}{What are the key metrics for success?}\label{faq:metrics}
  \begin{enumerate}[nosep]
    \item \textbf{External retention}: of developers who install Lux outside
      Punt Labs, the share still running it after one week. This is the metric
      that converts the internal-only evidence into external validation.
    \item \textbf{Companion-tool pull}: share of Lux installers who adopt at
      least one other Punt Labs tool (biff, vox, quarry). This measures the
      ecosystem thesis directly.
    \item \textbf{Element-coverage progress}: element kinds migrated onto the
      Element-ABC path (4 of 25 today), tracking the v2 buildout.
    \item \textbf{Interaction usage}: sessions that use the ask-the-user loop,
      distinguishing Lux-as-display from Lux-as-interaction-surface.
  \end{enumerate}

  Decision threshold: if external one-week retention stays near zero after a
  genuine external trial, the product is an internal tool, not a market one, and
  the roadmap should stop treating external adoption as a near-term goal.
\end{faqpair}

\begin{faqpair}{Why now? What has changed?}
  Three things converged. First, terminal-hosted agents, Claude Code and Codex,
  became the most capable agentic coding tools, and developers now run
  several at once, all producing output invisible beyond scrollback. Second, the
  vendor response is web dashboards that pull the agent onto remote compute,
  opening a clear gap for a local-first visual surface. Third, Lux itself proved
  the display and the interaction loop have real daily value internally, and the
  Hub/Display architecture that makes the shared, attributed, multi-agent window
  possible is shipped, not speculative.
\end{faqpair}

\begin{faqpair}{What are we explicitly not building?}\label{faq:wont}
  Lux is a curated visual surface, not a platform for arbitrary code. It will not
  become a self-extending display where agents author and register their own
  extensions; there is no extension registry, no community extension store, and
  no in-display LLM authoring agent. It exposes no HTTP or Ollama-shaped service
  API: MCP is the interface. It has one rendering mode (native ImGui), not
  several. It does not host local model inference, does not render to the web or
  Electron, and does not target mobile or tablet. Applications and applets
  (clocks, calculators, standalone viewers) are out of scope. These exclusions
  are what keep the built-in vocabulary coherent and the trust model simple; the
  full list is in the Feature Appendix.
\end{faqpair}

% ══════════════════════════════════════════════════════════════════════════════
% FOUR RISKS ASSESSMENT
% ══════════════════════════════════════════════════════════════════════════════

\newpage
\section*{Risk Assessment}

\renewcommand{\arraystretch}{1.6}
\begin{tabularx}{\textwidth}{@{} l l X @{}}
\toprule
\textbf{\color{SectionBlue}Risk} & \textbf{\color{SectionBlue}Rating} & \textbf{\color{SectionBlue}Assessment} \\
\midrule
Value      & \textcolor{RiskRed}{High} &
  All demonstrated value is \emph{internal}: the current users are one Punt Labs
  developer plus his agent fleet (see~\faqref{faq:evidence}), not the target
  customer the product is built for: external developers running several
  terminal agents. ``Runs in every session'' counts agent sessions, not distinct
  humans; external adoption is zero. Validation confined to the company that
  built Lux is a High Value-risk signal in the four-risks framework, so this row
  is rated High, not Medium. It is why \faqref{faq:next-step} makes the day-7
  external trial the single most important next step. \\
Usability  & \textcolor{RiskAmber}{Medium} &
  Install is one command and the display opens automatically on first output, so
  the intended onboarding is simple. But that flow is proven \emph{internally}
  only; it has never completed on a machine outside Punt Labs. Real prerequisites
  add friction: Python~3.13+, macOS/Linux only, and a curl-pipe-to-sh install
  that asks the user to trust a script. The day-7 external cohort
  (see~\faqref{faq:next-step}) is the first real measurement of
  onboarding-completion; until then the Low-friction claim is a design intent,
  not a measured result. \\
Feasibility & \textcolor{RiskAmber}{Medium} &
  The core is shipped and works: Hub/Display, MCP front door, the data-display
  core, the interaction loop, and a formally-verified singleton lifecycle. The
  open work is the v1-to-v2 element migration: 21 of 25 kinds still on the
  legacy path, with a handler-wrapping seam refactor and the table kind as the
  hard cases. No single step is high-risk, but the aggregate is substantial and
  is the primary feasibility exposure. Mitigated by incremental crossing with the
  legacy path deleted last (see~\faqref{faq:roadmap}). \\
Viability  & \textcolor{RiskAmber}{Medium} &
  MIT-licensed with no direct revenue. Strategic value depends entirely on
  ecosystem pull for the companion Punt Labs tools. If Lux does not drive
  adoption of biff, vox, and quarry, it is a public good without business
  return. Acceptable only if treated as infrastructure, which is the stated
  intent (see~\faqref{faq:revenue}). \\
\bottomrule
\end{tabularx}
\renewcommand{\arraystretch}{1.0}

% ══════════════════════════════════════════════════════════════════════════════
% FEATURE APPENDIX
% ══════════════════════════════════════════════════════════════════════════════

\newpage
\section*{Feature Appendix}

Defines the scope boundary for Lux. Must Do is the shipped v1 identity; Should
Do is the v2 direction; Won't Do is the firm exclusion set that keeps the
built-in vocabulary coherent and the trust model simple.

\subsection*{Must Do}

The core that makes Lux what it is. These ship in v1.

\begin{enumerate}[nosep,leftmargin=2.5em]
  \featureitem{Typed JSON element protocol}{agents describe UI as a tree of
    typed elements: tables, text, plots, groups, buttons, and more; the
    protocol is the API surface}\label{feat:protocol}
  \featureitem{Native GPU-accelerated rendering}{a persistent, shared ImGui
    window painting every frame on the developer's own GPU, with no web stack
    and no vendor compute}\label{feat:render}
  \featureitem{Hub/Display io-model}{\texttt{luxd} holds authoritative UI state;
    the display renders a replica; UI state crosses the wire, render calls do
    not}\label{feat:hubdisplay}
  \featureitem{Data-display core}{filterable, searchable tables with row
    selection and detail panels; dashboards; diagrams; images; 25 element
    kinds in all}\label{feat:data}
  \featureitem{Ask-the-user interaction loop}{an agent poses a question through a
    dialog, checkbox, or button and receives the answer back via Hub-side handler
    dispatch}\label{feat:interaction}
  \featureitem{Two-tier remote handler dispatch}{the Hub owns handlers; the
    Display forwards interactions back and the real handler runs on the Hub's
    copy}\label{feat:d21}
  \featureitem{Multi-agent shared window}{multiple agents share one window, each
    in its own frame with clear attribution}\label{feat:multiagent}
  \featureitem{MCP front door}{a first-class entry point where \texttt{mcp-proxy}
    bridges Claude Code stdio to the \texttt{luxd} WebSocket}\label{feat:mcp}
  \featureitem{Introspection and control surface}{\texttt{inspect\_scene},
    \texttt{list\_scenes}, \texttt{screenshot}, and related tools let an agent
    verify what actually rendered}\label{feat:introspect}
  \featureitem{One-command install with auto-open display}{a single curl command
    installs the package and plugin; the window opens on first
    output}\label{feat:install}
  \featureitem{Formally-verified singleton lifecycle}{the spawn/reap/bind
    concurrency (the display singleton lifecycle only, not rendering,
    protocol, or dispatch) is specified in Z and ProB-model-checked for
    singleton-serving, never-unlink-live, no-two-winners, clean-exit, and
    deadlock-freedom}\label{feat:formal}
\end{enumerate}

\subsection*{Should Do}

The v2 direction. Valuable, not v1-blocking.

\begin{enumerate}[nosep,leftmargin=2.5em]
  \featureitem{Complete ImGui element coverage}{every ImGui widget available as a
    curated, built-in element; the point is not coverage for its own sake but the
    input primitives that enable rich GUI-driven interaction with the
    agent}\label{feat:coverage}
  \featureitem{Element-ABC migration}{move the remaining 21 legacy element kinds
    onto the distributed Element-ABC / Hub-Display path, retiring the legacy
    render dispatch when the last kind crosses}\label{feat:migration}
  \featureitem{Texture-cache eviction policy}{bound the image texture cache,
    which today grows without eviction}\label{feat:texture}
  \featureitem{Companion-tool surfacing}{render biff presence panels, vox
    controls, and quarry search as first-class Lux views alongside the beads
    board}\label{feat:companion}
\end{enumerate}

\subsection*{Won't Do}

Explicitly excluded to keep the vocabulary curated and the trust model simple.

\begin{enumerate}[nosep,leftmargin=2.5em]
  \featureitem{Self-extending extension registry}{agents will not author and
    register arbitrary display extensions at runtime; the fixed, curated element
    set already composes without limit, so infinite element kinds are
    unnecessary, and Lux stays a built-in vocabulary rather than a plugin
    platform}\label{feat:no-registry}
  \featureitem{Community extension store}{no marketplace of third-party
    extensions with ratings and verification badges, which sit outside the trust
    model Lux commits to}\label{feat:no-store}
  \featureitem{In-display LLM authoring agent}{Lux does not embed a model that
    writes new UI code inside the display}\label{feat:no-agent}
  \featureitem{HTTP or Ollama-shaped service API}{MCP is the interface; Lux does
    not expose a general HTTP service surface}\label{feat:no-http}
  \featureitem{Multiple rendering modes}{one native ImGui mode only, with no
    texture-buffer or separate-OS-window rendering modes}\label{feat:no-modes}
  \featureitem{Applications and applets}{clocks, calculators, docks, and
    standalone viewers are out of scope}\label{feat:no-apps}
  \featureitem{Local model inference}{Lux renders; it does not host models}\label{feat:no-inference}
  \featureitem{Web or Electron rendering}{Lux renders natively via ImGui on the
    developer's GPU. HTML, Electron, and webviews are out of scope, since web
    rendering leads to the vendor-compute dependency Lux exists to
    avoid}\label{feat:no-web}
  \featureitem{Mobile or tablet}{desktop only (macOS, Linux)}\label{feat:no-mobile}
\end{enumerate}

% ══════════════════════════════════════════════════════════════════════════════
% REFERENCES
% ══════════════════════════════════════════════════════════════════════════════

\newpage
\printbibliography[title={References}]

\end{document}
